\section{Discussion}
\label{sec:discussion}

\subsection{Physical Significance of the \texorpdfstring{$K=3$}{K3} Tri-Modal Model}

The mathematical significance of the $K=3$ peak identified by the silhouette analysis is consistent with a deep physical pattern in stratified geophysics. Historically, stable boundary layer regimes have been modeled as binary states: either fully turbulent or weakly turbulent/laminar. This binary classification forms the basis of traditional patch-gated look-up tables used in mesoscale forecasting.

Our metric-consistent framework reveals that a binary assumption is physically incomplete. The stable boundary layer is a tri-modal system. The critical role of $K=3$ is to isolate the \emph{Intermittent Shear Bursting State} (Regime 3) as a long-lived, statistically independent physical regime rather than a fleeting numerical artifact. 

While the continuous turbulence mode (Regime 1) handles standard isotropic mixing and the wave-dominated mode (Regime 2) handles horizontal wave propagation, Regime 3 tracks the non-equilibrium energy exchange between those scales. This state represents periods where the vertical structure is highly non-monotonic, characterized by localized, suspended shear maxima that generate independent internal mixing layers aloft while remaining completely isolated from the surface. Physically, this intermittent burst signature is consistent with at least three well-documented sub-mesoscale phenomena in the nocturnal boundary layer: top-down shear instability driven by the nocturnal low-level jet (LLJ) decoupling from the surface layer; localized drainage flows and density currents impinging on the base of the stable inversion; and solitary-wave breaking events that locally steepen the potential temperature gradient until Kelvin--Helmholtz instability ensues \citep{Sun2015}. The elevated spectral curvature ($\chi_N$) associated with Regime~3 (Table~\ref{tab:gmm-centroids}) remains a useful roughness descriptor, and the regenerated feature correlation matrix indicates substantial negative coupling with $D_{\mathrm{eff}}$. We therefore interpret $(D_{\mathrm{eff}},\chi_N)$ as a coupled dual-coordinate pair rather than independent orthogonal axes: $D_{\mathrm{eff}}$ tracks modal occupancy breadth, while $\chi_N$ tracks weighted tail sharpness within that occupancy.
In formal terms, this transition behavior is captured by the non-orthogonal interaction block in Eq.~(\ref{eq:energy_overlap}), where off-diagonal terms $\langle \mathbf{c}^{(j)}, \mathbf{M}\mathbf{c}^{(k)} \rangle$ remain finite and encode cross-scale exchange rather than additive partition residuals.

The observed model-rank collapse by $\Delta_{\mathrm{rank}} = 24$ during these shifts is consistent with a reduced-order manifold neighborhood as an active physical characteristic of the nocturnal inversion layer, and tracking this state provides a promising framework for resolving missing sub-grid scale parameterizations.

\subsection{Observation, Interpretation, and Closure Hypotheses}

To keep the theoretical claims reviewer-safe, we separate evidence into three levels. Level 1 consists of direct diagnostics: reconstructed coordinates, cluster structure, and time-local covariance patterns. Level 2 is geometric interpretation: folded-sheet language, projection singularity, and the reduced-order reading of $D_{\mathrm{eff}}$ as a proxy for contraction and relaxation. Level 3 is closure hypothesis: whether these coordinates can be embedded in an autonomous or weakly forced slow manifold with Fenichel-style persistence. The present manuscript supports Levels 1 and 2; Level 3 remains a testable continuation program.

\subsection{Advancing Beyond MOST}

By projecting sparse meteorological tower observations onto a metric-consistent Riemannian manifold, this study addresses a long-standing diagnostic bottleneck in stable boundary layer meteorology. The application of our pseudospectral architecture to the CASES-99 intensive observation period shows that the traditional, highly variable continuum of nocturnal stable flows can be separated into three statistically distinct regimes without relying solely on empirical flux-profile look-up tables.

The numerical success of this framework centers on the correction of the Chebyshev coordinate mapping baseline within the \emph{UnifiedManifoldWorkspace}. Resolving the array-ordering inversion in the node distribution ensured that the observed instrument heights (\SI{1.5}{\meter} to \SI{55.0}{\meter}) were accurately mapped across the computational domain $\xi \in [-1, 1]$, preventing the structural rank collapse that previously forced metrics into an uninformative rank-1 subspace. In the current synoptic audit, campaign-window means of $D_{\mathrm{eff}}$ are partially collared across regimes and therefore are not treated as a stand-alone separator. Regime validity is instead established from the multivariate geometry of $[D_{\mathrm{eff}}, F_W, \chi_N]$, consistent silhouette stability, and independent SVD conditioning diagnostics ($r_{\mathrm{eff}}$ near the sensor-space ceiling with compressed wave-support subspaces). This interpretation supports the same physical conclusion while avoiding over-attribution to any single scalar metric.
This interpretation also preserves the agnosticism boundary established in Methods: state coordinates are transferable across observing geometries, but confidence intervals are singular-spectrum conditioned (Eqs.~(\ref{eq:svd_uncertainty_map})--(\ref{eq:svd_covariance})). Consequently, sparse vertical gaps or level clustering should be reported as uncertainty inflation mechanisms rather than treated as coordinate-system failure.

Furthermore, the regenerated feature correlation matrix (Table~\ref{tab:corr-matrix}) offers a path forward for overriding the current parameterization failures of Monin--Obukhov similarity theory (MOST). Because $\chi_N$ is strongly coupled to $D_{\mathrm{eff}}$ while $\mathrm{Ri}_b$ remains comparatively weakly coupled to the manifold diagnostics, $\chi_N$ should be treated as a curvature-weighted companion coordinate to modal entropy, not as a fully independent geometric axis. This framing still supports the physical interpretation of intermittent bursts, but avoids overcommitting to any single fixed correlation coefficient before final uncertainty bracketing. In this context, Interaction Residual extremes are most consistently interpreted as event-like wave--turbulence coupling signatures rather than instrumental or numerical contamination \citep{Sun2015, Mahrt2010}.
Classical MOST-style closures implicitly assume near-orthogonal scale separation, but the persistent off-diagonal interaction terms in Eq.~(\ref{eq:energy_overlap}) show that intermittent transition intervals violate that assumption. This provides a concrete mathematical reason why binary turbulence/wave closures underperform during nocturnal regime switching.

% ============================================================================
